% 08-privacy.tex: the privacy boundary, stated exactly, and measured where it was measured.
\section{The privacy boundary: what a verifier is allowed to learn}
\label{sec:privacy}

\motif{Current authorization is not privacy.}

\analogy{At some gates the guard learns only that the traveller belongs to the town, and not which townsperson is standing there; what the guard may write down is decided by the town, not by the guard.}

Two things in Polaris use zero knowledge, one thing deliberately does not exist, and the distinction
between them is the most important sentence in this section. Version 3 adds a fourth thing: an
adversary, written to falsify the claim, and what it found.

\subsection{Three disclosure levels}

Every verification event is recorded at one of three levels. \code{ZERO\_KNOWLEDGE} proves that a
valid token exists without recording which: the event's token identifier is null, and a check
constraint refuses any row that says otherwise (C2). \code{SELECTIVE} records named attributes only.
\code{FULL} records the identity and is logged and rate-limited. The level is decided by the server
and cannot be upgraded by the client (C6): a relying party cannot turn a zero-knowledge check into a
full one by asking differently, and the Atlas, the operators' bounded console (Section~\ref{sec:cognition}), redacts a zero-knowledge event's location server-side.
The default level is zero-knowledge. \sImpl

\subsection{Issuer-side unlinkable verification records}

The first use of zero knowledge is a property of the issuer's database. A default-mode verification
stores no token identifier, so the verification graph, who was verified where and when, cannot be
reconstructed from the database even by someone holding all of it. This is a structural claim, not a
promise: the repository carries a redaction proof that the redacted field is non-derivable under a
stated adversary model, property tests exercise every read path, and the claim that two relying
parties pooling everything they hold still cannot link a holder is a formal model checked on every
push. The property test that measures reconstruction risk runs an adversary against isolated
zero-knowledge events and asserts its success rate is no better than chance, and the assertion is
shaped so that a weak adversary cannot flatter the result. The claim is scoped precisely. It covers
zero-knowledge-mode events. Selective and full events do carry a token identifier, because that is
what those levels mean. \sImpl

\subsection{The membership proof}

\figfloat{holder-flow}{Holder, relying party and issuer, with what each sees and what each is
prevented from seeing. The holder presents a file; the relying party decides offline with the
detached verifier or an SDK and learns that a valid credential exists and its value; the issuer
records, in zero-knowledge mode, an event with no token identifier, and keeps no who-verified-whom
row for the online status check. The three flows form a circle turning one way; at its centre, in red, is the stated limitation: on a plain presentation the
credential value is a stable handle across the verifiers it is shown to.}{fig:holder-flow}

The second use of zero knowledge is a proof object. A Merkle tree behaves like a cryptographic
fingerprint of an entire set: instead of handing a verifier the whole set, one can hand it a short
proof that a particular item belongs to the set with the committed fingerprint. Polaris commits, per
epoch, a Merkle root over the active-token set, and a holder can produce a Plonky2 proof (a hash-based zero-knowledge proof system) on their own
device, against a published anonymity set, that their token is a leaf under that root, bound to the
epoch, a context and a nonce, without revealing which leaf. The verifier sees the proof bytes and the
public inputs and cannot recover the leaf; that is the proof system's property, not an API rule. The
setup is transparent, with no ceremony, and the commitments are hash-based. The default tree depth
of fourteen covers the schema's ten-thousand-leaf epoch cap, so the anonymity set is the epoch, and
the epoch has a floor: an epoch smaller than the minimum anonymity set (a database setting, twenty
unless set; the notional sample sets one and says so) is refused rather than closed smaller, and a proof against an older one below it answers \emph{privacy unavailable}
rather than verifying. The authority waits for members or lengthens its cadence; epochs are not
merged. The detached verifier applies the same floor offline to a foreign authority's epoch, from the
size its signed checkpoint states.
On the reference machine a proof takes about 25~ms to make and about 12~ms to check, and the verdict
is two-witnessed at the statement level. That second witness is independent in implementation
only: it is a separate program, written by the same author from the same reading of the circuit,
so it catches an implementation slip and not a misreading of the statement. The two signature
witnesses (Section~\ref{sec:signatures}) are independent in implementation and in authorship, and
share only the standard. Neither pair is independent in its interpretation of the specification.
The proof carries a per-verifier nullifier, a value that repeats if the same holder proves twice in the
same scope and reveals nothing else,
\code{Poseidon(secret\,||\,scope\,||\,epoch)}, Poseidon being a hash function designed for proof
circuits, so a relying party can enforce \emph{one enrolled
identity, once} (not one human: C3 does not see a human enrolled twice, Section~\ref{sec:core})
within its own scope while another cannot correlate its nullifiers. \sImpl

What the proof does not do is stated with the same care. It proves membership and binding, nothing
else: whether a token is active, permitted for the context and unrevoked is decided when the epoch is
composed, not inside the circuit. The nonce prevents substituting a captured proof into another
context; online, a single-use nonce store refuses a replayed bundle, and offline the verifier keeps no
state, so a relying party that needs single use issues a challenge. The nullifier does not hide the
holder from the \emph{issuer}, which derives every leaf in order to build the tree and could compute
any member's nullifier in any scope; the property is between relying parties. And the circuit over
the proof library has had no external cryptographic audit; the repository's soundness ledger says
not to protect real identities with it as shipped. The proof library has been evaluated against its
successor and kept, with the triggers that would reopen the question recorded.
\sExt{} for the audit.

\subsection{Bounded, not permanent, and what bounded was found to mean}

Polaris is a full-credential presentation system, and the consequence would be flat if the login
token's subject were the credential's hash: the same value at every relying party, so that two of
them could join their user tables exactly and forever. The subject and the presentation handle are
derived under the relying party's own scope, stable where an account needs it and
unrecognisable at the next verifier; since the subject is the value a relying party \emph{writes
down}, and written-down values are the ones pooled, sold, subpoenaed and breached, this is the half
that matters most in practice. The other half did not change: a plain presentation still shows the
verifier a stable \code{token\_value}, the issuer's signature and the holder's public key, so the
guarantee is about what a verifier should \emph{store}, not about what it is \emph{shown}, and the
verifier reports which it is giving: \code{exposed} for a plain presentation, \code{bounded} for the
zero-knowledge path. \sImpl

Version 3 can say what \code{bounded} was measured to mean, because under the operating contract the
repository's lab asked the question directly: given complete transcripts from two presentations,
what advantage does a colluding pair of verifiers have at deciding \emph{same holder}? Three findings,
in the order they were made.

\begin{description}
  \item[The verdict did not mean what it said.] \code{bounded} was reported whenever a scoped
    nullifier was present. It never checked the premise its own documentation stated. A presentation
    carrying the nullifier \emph{and} the credential's stable token value, issuer key and signature
    was reported bounded while two verifiers could link it with a single string comparison. Fixed the
    same day: \code{bounded} now requires the nullifier and the absence of every field that is
    identical across verifiers, including the holder's public key, which is stable by construction and
    is exactly why the pairwise handle is derived from it rather than shown. The fix corrected a
    verdict; it made no presentation more private than it was.
  \item[No advantage beyond the anonymity set, and nothing floors it.] An adversary that flattens
    every scalar in a transcript, drops what is constant across the population and matches on the
    rest, with a positive control that links the exposed population at an accuracy of one, achieves
    on the bounded population an accuracy that tracks $1/N$ at every size from two holders to two
    hundred. Within an epoch, scoping does what the word implies against this adversary. And
    $N$ is the epoch's membership. The schema constrains an epoch's committed count to between one
    and ten thousand and nothing floors it above one, so an epoch that closes with three members
    gives a colluding pair a one-in-three guess, and one that closes with a single member identifies
    the holder outright, while the verifier still, correctly, reports the transcript as bounded. The
    count is readable off the signed checkpoint; nothing requires a relying party to look, and no
    verdict mentions it. Reporting it would be a new product guarantee, which lab work does not get
    to create, so it is recorded as a limitation for a deploying organisation.
  \item[What the adversary cannot see.] It matches on exact equality and nothing else. Shown a
    population that is perfectly linkable with no equal field in it, it stays at chance. So the null
    result means no field is \emph{identical} across verifiers, not that none is \emph{correlated}.
    The one channel it can read is transcript length, which carried the entire above-chance residue
    in that control; nothing in a bounded Polaris presentation varies in size per holder today, and a
    deployment where disclosed attribute sets or optional elements made it vary would open it.
    Presentation size, timing, issuer metadata, status artifacts and transcript structure are
    unmeasured, and the lab says so rather than extrapolating.
\end{description}

The path a stranger's wallet actually used cuts the other way. An SD-JWT VC presentation over
OpenID4VP, measured the same way, leaves nine values identical across two verifiers, the holder's
key, the issuer's signature, the disclosure digests, the hash of the presentation and its length,
all issuer-signed so the holder cannot vary them. That is a known property of presenting one such
credential repeatedly; the ecosystem's mitigation is batch issuance, many credentials with one holder
key each, which the repository does not implement. Adopting the ecosystem's format imported the
ecosystem's linkability; it was not visible before it was measured, and it is not a reason to stop
interoperating. So cross-verifier correlation is bounded on the native zero-knowledge path,
measured within the limits above, and exposed on the interoperable one. \emph{Unlinkable by
default} still means the issuer's zero-knowledge-mode records. \sImpl{} for the derivations and the
measurement, with the residue stated rather than closed.

\subsection{Bounded surfaces}

The remaining privacy mechanisms bound what a surface can return rather than what is stored. Every
map and API aggregate is hard-capped (C8), so the whole population cannot be extracted through the
operator's atlas, and a request for six and a half million cluster bins returns at most five
thousand. The single-entity investigation pages are built to refuse cross-individual link analysis
and are audited: a meta-audit table records who read which audit surface, with what filter, and how
many rows came back, and it is append-only. The authority's most invasive power, pulling one named
person's whole verification history under a warrant, writes the same row, and a structural rule
extends that to any read hidden behind a procedure or a view, where a search for a \code{SELECT}
would not find it. The authority-structure layer, Athena, is forbidden by
\code{check\_athena\_no\_person} from referencing any natural-person table or column (Section~\ref{sec:cognition}). And
timestamps, the timestamp authority and the exchange gateway are built to learn nothing about what
they touch, which Section~\ref{sec:institutions} describes. \sImpl

\begin{layercard}{Layer card: the privacy boundary}
\qa{What it is}{Three disclosure levels with a CHECK behind the default; server-side enforcement; the redaction proof and its adversary; the epoch membership proof, its nullifier and its second witness; the per-verifier handle; bounded aggregates; an audited single-entity investigation surface.}
\qa{Why it exists}{A verification log is a surveillance backbone unless it is structurally unable to be one.}
\qa{What it trusts}{The C2 constraint; the proof library's soundness; the redaction proof's adversary model, and the linkability adversary's stated blind spots.}
\qa{What it refuses}{A client's claim about its own disclosure level; a token identifier on a zero-knowledge row; an unbounded read; a person node in the authority layer; a verdict of \code{bounded} over a transcript with any field identical across verifiers.}
\qa{What it can see}{At zero knowledge: that a valid token was verified at a time in a context. At full: the identity, logged.}
\qa{What it cannot see}{Which token, in zero-knowledge mode, even from the whole database; and the size of the crowd a given proof hides in, which no verdict reports.}
\qa{If it fails}{An identifier on a zero-knowledge row is unstorable; a cap removed fails the build; a person reference in Athena fails \code{check\_athena\_no\_person}; an adversary weakened to flatter the result fails the property suite's own control. Cross-verifier correlation on a plain presentation is not a failure; it is the documented boundary.}
\qa{Now or later}{Now. Linkability-resistant presentation for the interoperable path, and a crowd-size floor, are a different system and later.}
\qa{Inside and outside}{Inside: status. Outside: the holders and relying parties who use the credential.}
\end{layercard}

\nextlayer{A credential bounded in what it reveals still has to be used: held, presented, accepted
and logged in with. The next layer is where the credential meets software that is not the issuer's,
and where, for the first time, that software was not the author's either.}
